% ══════════════════════════════════════════════════════════════════
%  FICHE D'EXERCICES — Chapitre 2 : Les alcanes
% ══════════════════════════════════════════════════════════════════
\section{Exercices type Baccalauréat}
\textit{Données : $M(\ce{H})=\SI{1}{\gram\per\mol}$, $M(\ce{C})=\SI{12}{\gram\per\mol}$ ;
$V_m=\SI{24}{\litre\per\mol}$ ; densité d'un gaz par rapport à l'air : $d=M/29$.}

\rubrique{Nomenclature \& isomérie}

\exo{}\diff{1} Nommer les alcanes suivants :
a)~$\ce{CH3-CH(CH3)-CH2-CH2-CH3}$ ;
b)~$\ce{(CH3)3C-CH2-CH3}$ ;
c)~$\ce{CH3-CH(C2H5)-CH(CH3)-CH3}$.

\exo{}\diff{1} Écrire les formules semi-développées de : a)~2,2,4-triméthylpentane ;
b)~3-éthyl-2-méthylhexane ; c)~4-éthyl-3,3-diméthyloctane.

\exo{}\diff{2} Déterminer le nombre d'isomères de constitution de l'hexane $\ce{C6H14}$,
écrire toutes leurs formules semi-développées et les nommer.

\rubrique{Combustion \& détermination de formule}

\exo{}\diff{1} Écrire et équilibrer l'équation de la combustion complète du propane
$\ce{C3H8}$, puis du butane $\ce{C4H10}$.

\exo{}\diff{2} La combustion complète de \SI{2.9}{\gram} d'un alcane gazeux $A$
nécessite \SI{7.2}{\litre} de dioxygène (\si{\litre} mesurés à $V_m=\SI{24}{\litre\per\mol}$).
\begin{enumerate}[label=\arabic*)]
\item Établir l'équation générale de combustion d'un alcane $\ce{C_nH_{2n+2}}$.
\item Déterminer la formule brute de $A$, puis ses isomères et leurs noms.
\end{enumerate}

\exo{}\diff{2} Un alcane $A$ a une densité de vapeur par rapport à l'air égale à
$d = 2{,}48$. Déterminer sa masse molaire, sa formule brute, et le nombre de ses
isomères.

\exo{}\diff{3} Un mélange de méthane $\ce{CH4}$ et d'éthane $\ce{C2H6}$ a un volume
total de \SI{2.0}{\litre}. Sa combustion complète produit \SI{3.0}{\litre} de
dioxyde de carbone (volumes mesurés dans les mêmes conditions).
\begin{enumerate}[label=\arabic*)]
\item Écrire les équations de combustion des deux gaz.
\item Déterminer les volumes de méthane et d'éthane dans le mélange.
\item En déduire le volume de dioxygène consommé et la composition molaire (\%) du mélange.
\end{enumerate}

\rubrique{Réactivité : substitution radicalaire, craquage}

\exo{}\diff{1} Écrire l'équation de la monochloration du méthane. Préciser les
conditions de la réaction et nommer le produit organique obtenu.

\exo{}\diff{2} La chloration du 2-méthylpropane $\ce{(CH3)3CH}$ peut conduire à deux
produits monochlorés isomères.
\begin{enumerate}[label=\arabic*)]
\item Écrire leurs formules semi-développées et les nommer.
\item Détailler le mécanisme radicalaire (initiation, propagation, terminaison)
pour la formation de l'un d'eux.
\end{enumerate}

\exo{}\diff{2} Le craquage d'une molécule de décane $\ce{C10H22}$ peut produire un
alcane et un alcène. Écrire deux équations de craquage possibles et nommer les
produits.

\exo{}\diff{2} On réalise la chloration complète (substitution de tous les H) du
méthane. Écrire les équations successives et nommer les quatre produits chlorés obtenus.

\rubrique{Problèmes de synthèse — type Baccalauréat}

\sujetbac[d'après Baccalauréat D, Cameroun]
Un hydrocarbure saturé gazeux $A$ a une masse molaire $M = \SI{72}{\gram\per\mol}$.
\begin{enumerate}[label=\arabic*)]
\item Montrer que $A$ est un alcane et déterminer sa formule brute.
\item Donner les formules semi-développées et les noms de tous les isomères de $A$.
\item L'un de ces isomères ne donne qu'un seul produit de monochloration.
L'identifier, le nommer, et écrire l'équation de sa monochloration.
\item Calculer le volume d'air ($\SI{20}{\percent}$ de \ce{O2} en volume) nécessaire
à la combustion complète de \SI{3.6}{\gram} de $A$.
\end{enumerate}

\sujetbac[Analyse d'un gaz domestique]
Le gaz d'un réchaud est un mélange de propane $\ce{C3H8}$ et de butane $\ce{C4H10}$.
La combustion complète de \SI{1}{\mol} de ce mélange consomme \SI{6.3}{\mol} de
dioxygène.
\begin{enumerate}[label=\arabic*)]
\item Écrire les équations de combustion des deux alcanes.
\item Déterminer la composition molaire (\%) du mélange.
\item En déduire la masse molaire moyenne du mélange et la masse de \ce{CO2} produite
par la combustion de \SI{1}{\mol} de mélange.
\end{enumerate}
